% GENERATED FILE. Edit canonical lessons, then run npm run generate:books.
% canonical-source-hash: fnv1a64:8b1828499843e1ea
% canonical-lessons: MR-C304-varnan-karne, MR-C304-prasang, MR-C304-ulgadne, MR-C304-gajne, MR-C304-pasarne

\chapter{Telling It Well}
\label{ch:mr-telling-it-well}
\hlchaptermodality{304}

\begin{chapteropening}
\textbf{By the end of this chapter:} \emph{I can describe an episode and say how the news spread.}

Complete the last lesson of chapter 304: I can describe an episode and say how the news spread.
\end{chapteropening}

\section[\texorpdfstring{\mr{वर्णन} \mr{करणे} (varṇan karṇe)}{वर्णन करणे (varṇan karṇe)}]{\mr{वर्णन} \mr{करणे} (varṇan karṇe) — to describe}

\label{lesson:MR-C304-varnan-karne}



\begin{quote}
Before the new one: say the Marathi for an anecdote, then the Marathi for to narrate.
\end{quote}

\subsection*{You'll want to know: \mr{वर्णन} \mr{करणे}}
\textbf{\mr{वर्णन} \mr{करणे}} — \emph{varṇan karṇe} — \textquotedblleft{}to describe\textquotedblright{}.

\begin{quote}
\mr{तिने} \mr{सहलीचं} \mr{सुंदर} \mr{वर्णन} \mr{केलं}.
\end{quote}

\begin{quote}
\emph{tine sahlīcaṁ sundar varṇan kelaṁ.}
\end{quote}

\begin{quote}
She described the trip beautifully.
\end{quote}

Past, with -\mr{ने} (he …): \mr{त्याने} \mr{वर्णन} \mr{केलं}, \emph{tyāne varṇan kelaṁ}.

What is described takes -\mr{चं}: \mr{गावाचं} \mr{वर्णन} \mr{करणे}.

\begin{cousinweb}
\textbf{\mr{वर्णन} \mr{करणे}} (\emph{varṇan karṇe}) joins \emph{varṇan}, 'a description', a tatsama taken from Sanskrit \emph{varṇana}, 'depicting, describing' (Monier-Williams), to \emph{karṇe}, 'to do'.
\end{cousinweb}

\begin{grammarlens}[title={What you've built}]
One more word for telling what happened.
\end{grammarlens}

\subsection*{Your turn}
\begin{itemize}
\raggedright
  \item \emph{Say it:} \emph{varṇan karṇe}
  \item \emph{Say it:} \emph{varṇan karṇe}, once more
  \item \emph{Say it:} \emph{tine sahlīcaṁ sundar varṇan kelaṁ.} — She described the trip beautifully.
  \item \emph{Recall:} say the Marathi for an anecdote, then the Marathi for to narrate, then say \emph{varṇan karṇe} again
\end{itemize}

\begin{tcolorbox}[breakable,colback=teal!4,colframe=teal!35!black,title={Before you move on}]
What does \mr{वर्णन} \mr{करणे} mean? (\textquotedblleft{}To describe\textquotedblright{}.) Say it once more.
\end{tcolorbox}
\section[\texorpdfstring{\mr{प्रसंग} (prasaṅg)}{प्रसंग (prasaṅg)}]{\mr{प्रसंग} (prasaṅg) — an occasion; an episode}

\label{lesson:MR-C304-prasang}



\begin{quote}
Before the new one: say the Marathi for to narrate, then the Marathi for to describe.
\end{quote}

\subsection*{You'll want to know: \mr{प्रसंग}}
\textbf{\mr{प्रसंग}} — \emph{prasaṅg} — \textquotedblleft{}an occasion; an episode\textquotedblright{}.

\begin{quote}
\mr{शाळेतला} \mr{तो} \mr{प्रसंग} \mr{मी} \mr{कधीच} \mr{विसरणार} \mr{नाही}.
\end{quote}

\begin{quote}
\emph{śāḷetlā to prasaṅg mī kadhīc visarṇār nāhī.}
\end{quote}

\begin{quote}
I will never forget that episode at school.
\end{quote}

Masculine.

\mr{प्रसंगी}, on occasion.

\begin{cousinweb}
\textbf{\mr{प्रसंग}} (\emph{prasaṅg}) is a tatsama, taken from Sanskrit \emph{prasaṅga}, 'attachment; an occasion, a juncture' (Monier-Williams).
\end{cousinweb}

\begin{grammarlens}[title={What you've built}]
One more word for telling what happened.
\end{grammarlens}

\subsection*{Your turn}
\begin{itemize}
\raggedright
  \item \emph{Say it:} \emph{prasaṅg}
  \item \emph{Say it:} \emph{prasaṅg}, once more
  \item \emph{Say it:} \emph{śāḷetlā to prasaṅg mī kadhīc visarṇār nāhī.} — I will never forget that episode at school.
  \item \emph{Recall:} say the Marathi for to narrate, then the Marathi for to describe, then say \emph{prasaṅg} again
\end{itemize}

\begin{tcolorbox}[breakable,colback=teal!4,colframe=teal!35!black,title={Before you move on}]
What does \mr{प्रसंग} mean? (\textquotedblleft{}An occasion; an episode\textquotedblright{}.) Say it once more.
\end{tcolorbox}
\section[\texorpdfstring{\mr{उलगडणे} (ulgaḍṇe)}{उलगडणे (ulgaḍṇe)}]{\mr{उलगडणे} (ulgaḍṇe) — to unfold; to come clear}

\label{lesson:MR-C304-ulgadne}



\begin{quote}
Before the new one: say the Marathi for to describe, then the Marathi for an occasion.
\end{quote}

\subsection*{You'll want to know: \mr{उलगडणे}}
\textbf{\mr{उलगडणे}} — \emph{ulgaḍṇe} — \textquotedblleft{}to unfold; to come clear\textquotedblright{}.

\begin{quote}
\mr{हळूहळू} \mr{त्या} \mr{घटनेचा} \mr{अर्थ} \mr{उलगडला}.
\end{quote}

\begin{quote}
\emph{haḷūhaḷū tyā ghaṭnecā artha ulgaḍlā.}
\end{quote}

\begin{quote}
Little by little the meaning of that event came clear.
\end{quote}

Past, with -\mr{ने} (he …): \mr{त्याने} \mr{उलगडलं}, \emph{tyāne ulgaḍlaṁ}.

Intransitive (\mr{अर्थ} \mr{उलगडतो}, the meaning unfolds) and transitive (\mr{कागद} \mr{उलगडणे}, to unfold a paper).

\begin{cousinweb}
\textbf{\mr{उलगडणे}} (\emph{ulgaḍṇe}) is inherited: CDIAL lists Marathi \emph{algaḍṇẽ}, 'to get unravelled; to unravel', most likely a variant of this verb, under \emph{*ullagyati}, 'is unfastened' (CDIAL 2363).
\end{cousinweb}

\begin{grammarlens}[title={What you've built}]
One more word for telling what happened.
\end{grammarlens}

\subsection*{Your turn}
\begin{itemize}
\raggedright
  \item \emph{Say it:} \emph{ulgaḍṇe}
  \item \emph{Say it:} \emph{ulgaḍṇe}, once more
  \item \emph{Say it:} \emph{haḷūhaḷū tyā ghaṭnecā artha ulgaḍlā.} — Little by little the meaning of that event came clear.
  \item \emph{Recall:} say the Marathi for to describe, then the Marathi for an occasion, then say \emph{ulgaḍṇe} again
\end{itemize}

\begin{tcolorbox}[breakable,colback=teal!4,colframe=teal!35!black,title={Before you move on}]
What does \mr{उलगडणे} mean? (\textquotedblleft{}To unfold; to come clear\textquotedblright{}.) Say it once more.
\end{tcolorbox}
\section[\texorpdfstring{\mr{गाजणे} (gājṇe)}{गाजणे (gājṇe)}]{\mr{गाजणे} (gājṇe) — to make a stir}

\label{lesson:MR-C304-gajne}



\begin{quote}
Before the new one: say the Marathi for an occasion, then the Marathi for to unfold.
\end{quote}

\subsection*{You'll want to know: \mr{गाजणे}}
\textbf{\mr{गाजणे}} — \emph{gājṇe} — \textquotedblleft{}to make a stir\textquotedblright{}.

\begin{quote}
\mr{तो} \mr{चित्रपट} \mr{त्या} \mr{वर्षी} \mr{खूप} \mr{गाजला}.
\end{quote}

\begin{quote}
\emph{to citrapaṭ tyā varṣī khūp gājlā.}
\end{quote}

\begin{quote}
That film made a big stir that year.
\end{quote}

Past (he …): \mr{तो} \mr{गाजला}, \emph{to gājlā}.

Intransitive; \mr{गाजलेला}, much talked about.

\begin{cousinweb}
\textbf{\mr{गाजणे}} (\emph{gājṇe}) is a tadbhava: CDIAL lists Marathi \emph{gājṇe}, 'thunders', under Sanskrit \emph{garjati}, 'thunders, roars' (CDIAL 4046). What thunders is heard everywhere.
\end{cousinweb}

\begin{grammarlens}[title={What you've built}]
One more word for telling what happened.
\end{grammarlens}

\subsection*{Your turn}
\begin{itemize}
\raggedright
  \item \emph{Say it:} \emph{gājṇe}
  \item \emph{Say it:} \emph{gājṇe}, once more
  \item \emph{Say it:} \emph{to citrapaṭ tyā varṣī khūp gājlā.} — That film made a big stir that year.
  \item \emph{Recall:} say the Marathi for an occasion, then the Marathi for to unfold, then say \emph{gājṇe} again
\end{itemize}

\begin{tcolorbox}[breakable,colback=teal!4,colframe=teal!35!black,title={Before you move on}]
What does \mr{गाजणे} mean? (\textquotedblleft{}To make a stir\textquotedblright{}.) Say it once more.
\end{tcolorbox}
\section[\texorpdfstring{\mr{पसरणे} (pasarṇe)}{पसरणे (pasarṇe)}]{\mr{पसरणे} (pasarṇe) — to spread}

\label{lesson:MR-C304-pasarne}



\begin{quote}
Before the new one: say the Marathi for to unfold, then the Marathi for to make a stir.
\end{quote}

\subsection*{You'll want to know: \mr{पसरणे}}
\textbf{\mr{पसरणे}} — \emph{pasarṇe} — \textquotedblleft{}to spread\textquotedblright{}.

\begin{quote}
\mr{ही} \mr{बातमी} \mr{संपूर्ण} \mr{गावात} \mr{पसरली}.
\end{quote}

\begin{quote}
\emph{hī bātmī sampūrṇa gāvāt pasarlī.}
\end{quote}

\begin{quote}
This news spread through the whole village.
\end{quote}

Past (he …): \mr{तो} \mr{पसरला}, \emph{to pasarlā}.

Intransitive; its transitive partner is \mr{पसरवणे}, to spread something.

\begin{cousinweb}
\textbf{\mr{पसरणे}} (\emph{pasarṇe}) is a tadbhava: CDIAL lists Marathi \emph{pasarṇe}, 'to spread out', under Sanskrit \emph{prasarati}, 'advances, extends' (CDIAL 8825).
\end{cousinweb}

\begin{grammarlens}[title={What you've built}]
One more word for telling what happened.
\end{grammarlens}

\subsection*{Your turn}
\begin{itemize}
\raggedright
  \item \emph{Say it:} \emph{pasarṇe}
  \item \emph{Say it:} \emph{pasarṇe}, once more
  \item \emph{Say it:} \emph{hī bātmī sampūrṇa gāvāt pasarlī.} — This news spread through the whole village.
  \item \emph{Recall:} say the Marathi for to unfold, then the Marathi for to make a stir, then say \emph{pasarṇe} again
\end{itemize}

\begin{tcolorbox}[breakable,colback=teal!4,colframe=teal!35!black,title={Before you move on}]
What does \mr{पसरणे} mean? (\textquotedblleft{}To spread\textquotedblright{}.) Say it once more.
\end{tcolorbox}
